\documentclass[12pt,a4paper,oneside,header=plain,tocpage=prefix]{report}
%----SELEÇÃO DE FONTES ----%
%\usepackage{mathpazo}%Palatino LinoType
%\usepackage[default]{cantarell}%Cantarell
\usepackage[garamond]{mathdesign}%Adobe Garamond Pro
%\usepackage[urw-garamond]{mathdesign}
%\usepackage{fourier}%Adobe Utopia
%\usepackage{mathptmx}%Times
%\usepackage{literat}%Literaturnaya
%\usepackage{newcent}  % New Century Schoolbook + Avant Garde + Courier
%\usepackage{bookman}  % Bookman + Avant Garde + Courier
%\usepackage{antt}  % Bookman + Avant Garde + Courier
%\usepackage[dayroman]{mathdesign}%Adobe Garamond Pro
%--------------------------%
\newcommand{\titulo}{Título da Monografia}
\newcommand{\autor}{Marciel Barros Pereira}
\newcommand{\orientador}{Prof. Dr. Sérgio Antenor de Carvalho}
\newcommand{\coorientador}{Prof. }
\newcommand{\membroa}{Prof. Dr. Nome do Professor Banca 1}
\newcommand{\membrob}{Prof. Dr. Nome do Professor Banca 2}
\newcommand{\membroc}{Prof. Dr Nome do Professor Banca 3, se existir}
\newcommand{\Local}{Fortaleza, Ceará}
\newcommand{\Data}{\today}
\newcommand{\DataDefesa}{\today}
%\newcommand{\DataDefesa}{XX de Janeiro de 2013}

\newcommand{\instituicao}{Universidade Federal do Ceará}
\newcommand{\departamento}{Departamento de Eng. de Teleinformática}
\newcommand{\curso}{Curso de de Graduação em Engenharia de Teleinformática}
\newcommand{\comentario}{Projeto de Final de Curso apresentado à Coordenação do {\curso} da {\instituicao} como requisito parcial para a obtenção do diploma de Engenheiro de Teleinformática.}
\usepackage{preambulo}
\usepackage{formatacao_marciel}%formatação pessoal - margens abnt, texto do capitulo em small caps (modificável)
\usepackage{lipsum}%gerador de texto

\usepackage[usenames,dvipsnames]{pstricks}
%\usepackage{epsfig}
\usepackage{pst-grad} % For gradients
\usepackage{pst-plot} % For axes
\begin{document}
\onehalfspacing
% ----------------------------------------------------------------------- %
% Onde serão inseridas informações que irão aparecer na capa e na
% folha de rosto
%
% Arquivo: capa.tex
% ----------------------------------------------------------------------- %
%\capa
%%capa feita manualmente
%%-------------------------
\begin{titlepage}
\begin{center}
\includegraphics[scale=0.33]{fig/logo_UFC.eps}\\
\vfill
{\MakeUppercase{\instituicao}}\par
{\MakeUppercase{\departamento}}\par
{\MakeUppercase{\curso}}\\
\vfill
\begin{center}
\Large\MakeUppercase{\titulo}\par
\end{center}
\vfill\vfill
\begin{center}
{\MakeUppercase{\autor}}
\end{center}
\vfill\vfill\vfill

\setlength{\parskip}{.3cm}
{\normalfont{\Local}} \par
{\normalfont{\Data}}
\end{center}
\end{titlepage}
%%-------------------------
%Folha de rosto feita manualmente
%----------------------------------------
\thispagestyle{empty}
\vspace{1.1cm}
  \begin{center}
    \MakeUppercase{\autor}
  \end{center}
\vfill\vfill\vfill
   \begin{center}
     \Large\MakeUppercase{\titulo}\par
   \end{center}
   \vspace{.8cm}
   \hspace{.45\textwidth}
     \begin{minipage}{.5\textwidth}
         {\comentario}\par
     \end{minipage}
\vspace{.8cm}
\begin{center}
   Orientador: {\orientador}
   \vspace{0.7cm}
   \end{center}
\vfill
\begin{center}
  \setlength{\parskip}{.3cm}
     \setlength{\parskip}{0cm}
     {\MakeUppercase{\instituicao}}\par
	{\MakeUppercase{\departamento}}\par
	{\MakeUppercase{\curso}}\\
     \setlength{\parskip}{.3cm}\par
\end{center}
\vfill\vfill
\begin{center}
\setlength{\parskip}{.3cm}
{\normalfont{\Local}} \par
{\normalfont{\Data}}
\end{center}

%---------------------------------------%arquivo de capa e folha de rosto
\thispagestyle{empty}%

\begin{center}
    %\includegraphics[width=2.5cm]{figs/UFC.eps} \\%
    \textsc{ \autor } \\
     \vspace{.5 cm} \textbf{ \titulo }     \\
\end{center}
    \vspace{.2 cm}
    Esta Monografia foi julgada adequada para a obten��o do diploma de Engenheira do Curso de Gradua��o em Engenharia de Teleinform�tica da Universidade Federal do Cear�.
    \assinatura{\autor}
    \vspace{0.2 cm}
     Banca Examinadora:
     \assinatura{\orientador \\ Orientador}
     \assinatura{\membroa\\}
     \assinatura{\membrob \\}
    \assinatura{\membroc\\}
%     \assinatura{\membrod}
     \vspace{0.2 cm}%\null\vfill%%

\begin{center}
    {\normalsize    Fortaleza, \today}
\end{center}

% ----------------------------------------------------------------------- %
% Pequena dedicatória ou uma epígrafe (uma citação pertinente ao seu 
% trabalho ou que represente o seu modo de pensar.) 
% 
%
% Arquivo: dedicatoria.tex
% ----------------------------------------------------------------------- %
\thispagestyle{empty}
\vspace*{\fill}

{ \raggedleft


\textit{Dedico este trabalho a ...}

~
}

% ----------------------------------------------------------------------- %
% Pequena dedicatória ou uma epígrafe (uma citação pertinente ao seu 
% trabalho ou que represente o seu modo de pensar.) 
% 
%
% Arquivo: dedicatoria.tex
% ----------------------------------------------------------------------- %
\thispagestyle{empty}
\vspace*{\fill}

{ \raggedleft


\textit{Epigrafe aqui...}

~
}

\chapter{Agradecimentos}
\label{CHP:ACKNOWLEDGMENT}%%
\thispagestyle{empty}
% \PARstartOne{A}{grade\c{c}o}
%Agrade\c{c}o primeiramente ao Pai Celestial por todas as ben\c{c}\~{a}os derramadas sobre mim
%durante toda minha vida.
%
%\`{A} minha esposa que est\'{a} sempre me apoiando desde a escrita da minha Monografia na gradua\c{c}\~{a}o.
%
%Ao Professor Dr. Paulo C\'{e}sar Cortez, meu Orientador Cient\'{\i}fico, pela amizade e a oportunidade de receber um pouco de sua sabedoria, pela disponibilidade
%apresentada e pelas condi\c{c}\~{o}es que me proporcionou para a realiza\c{c}\~{a}o
%deste trabalho.
%
%Ao Rodrigo Costa, pela paci\^{e}ncia e esfor\c{c}o compartilhado auxiliando na orienta\c{c}\~{a}o desse trabalho, e aos demais integrantes do grupo GIHM (Grupo de Inova\c{c}\~{a}o ).
%
%Aos amigos, Auzuir Ripardo, Pedro Pedrosa , Tarique Cavalcante, dentre
%outros integrantes do sub-grupo de pesquisa em Engenharia Biom\'{e}dica do Laborat\'{o}rio de Engenharia de Sistemas de Computa\c{c}\~{a}o, pela amizade e esfor\c{c}o compartilhado durante o curso de mestrado.
%
%Aos engenheiros de desenvolvimento do SIDI (Samsung Instituto de Desenvolvimento para a Inform\'{a}tica), Giovanni, Miguel e Nelson pelo embarque dos algoritmos na plataforma de testes.
%
%Por fim, mas n\~{a}o menos importante, \`{a} Coordena\c{c}\~{a}o de Aperfei\c{c}oamento de Pessoal de N\'{\i}vel Superior (CAPES) e Funda\c{c}\~{a}o Cearense de Apoio ao Desenvolvimento Cient\'{\i}fico e Tecnol\'{o}gico (FUNCAP) pelo suporte financeiro atrav\'{e}s do concedimento da bolsa de Mestrado.


% ----------------------------------------------------------------------- %
% Pequeno texto que em poucas palavras consegue expressar o trabalho.
% O resumo deve ser concebido de forma tal que, uma pessoa ao ler o resumo
% possa entender sobre qual assunto este trabalho trata.
%
% Arquivo: resumo.tex
% ----------------------------------------------------------------------- %
\pdfbookmark[1]{Resumo}{CHP:RESUM0}
\chapter*{Resumo}
\label{CHP:RESUM0}
\thispagestyle{empty}
\singlespacing
	\lipsum[3-4]
\onehalfspacing
\pdfbookmark[1]{Abstract}{CHP:ABSTRACT}
\chapter*{Abstract}
\label{CHP:ABSTRACT}%%
\thispagestyle{empty}


\PARstartOne{A} great number of web applications has, for mobile platforms \emph{iOS} and \emph{Android}, some version for mobile systems, which same information is shared on cloud. Some popular social networks like \emph{Twitter} and \emph{Foursquare} has communication interfaces REST as a service to other applications. This paper aims to make a comparative study between iOS and Android applications that use a RESTful web application. To conduct the study, it created a small social network to create quizzes using \emph{Ruby on Rails} and a SaaS service like \emph{Heroku}.


%\PARstartOne{H}{ard} disks are very important elements in any modern computer system and it is advisable to monitor the health of these devices because they are used for applications from personal computers to risk activities. This work aims to build a framework in Linux system using the standard commands \ac{ATA} and \ac{SCSI} to support the development of algorithms for test drives. To demonstrate the use of the framework, algorithms are implemented some test that is compared with diagnostic tools on the market. Finally, these algorithms are embedded in a bootable version of Linux, which allows you to perform during any computers that support booting from a live CD or USB stick, regardless of operating system installed. The test results show the algorithm is very satisfactory, demonstrating the effectiveness of the framework.

%For the development of the work all commands have been developed in ANSI C + + language. Several commands \ac{ATA} and \ac{SCSI} were implemented to support reading, self-testing and other features. Also, the algorithm are implemented from them.
%\newline
%\newline
\noindent \textbf{Keywords: iOS, Android, Mobile, Ruby on Rails, REST, RESTFul}.

%Hard disks are  elements in any modern computer system and it is  to monitor the health of these devices, because they are used for applications from personal computer to risk activities. This work aims to build a framework in Linux system using the standard commands ATA and SCSI to support the development of algoritms for test devices. 
% listas automáticas: sumário, lista de figuras e lista de tabelas
\pdfbookmark[1]{Sumário}{sumario_label} %\label{sumario_label}
\tableofcontents
\newpage
%---------------------------------------
\listoffigures
\addcontentsline{toc}{chapter}{Lista de Figuras}%Adiciona lista de figuras ao sumário
\newpage
%---------------------------------------
\listoftables
\addcontentsline{toc}{chapter}{Lista de Tabelas}%Adiciona lista de tabelas ao sumário
\newpage
%%Ajuste de espaçamento 1,5 linha ao restante do texto:
\onehalfspacing
%INCLUIR AQUI OS CAPÍTULOS
%
\chapter{Introdução}

%=======%
%=======
citacao 2 \cite{balanis3}
citacao 3 \cite{microwaverf}
citacao 4 \cite{arraybasics}
citacao 5 \cite{pozar4}
citacao 6 \cite{rfcircuit}
citacao 7 \cite{microwavehandbook}

\bibstyle{abnt-alf}
\bibliography{bibliografia}%Adicionar os 
\end{document}